\section{Graded Summary of Catholic Teaching on Hell}

The table below distinguishes the theological authority of the principal teachings and questions treated in the body. \textit{De fide} denotes truth formally revealed by God and proposed by the Church's solemn or ordinary universal Magisterium as requiring the assent of divine and Catholic faith; \textit{Sententia communis} denotes the steady consensus of the Fathers and Catholic theologians; \textit{Open theological question} denotes a matter on which revered Catholic teachers hold differing views within the bounds of orthodoxy.

\begin{threestable}{Teaching}{Grade}{Principal witnesses}
Existence of hell as the state of final damnation for demons and impenitent sinners & \textit{De fide} & Matt 10:28; 25:41, 46; \textit{Quicumque} (DS 76); Lateran IV (DS 801); \textit{Benedictus Deus} (DS 1002); CCC 1035 \\
Eternity of hell's punishment; rejection of temporary chastisement and universal restoration (\textit{apokatastasis}) & \textit{De fide} & Matt 25:46; Mark 9:43--47; 2 Thess 1:9; Synod of Constantinople (543, DS 411); Lateran IV (DS 801); CCC 1035 \\
Immediate descent into hell of souls dying in actual mortal sin & \textit{De fide} & Luke 16:22--23; Lyons II (DS 858); \textit{Benedictus Deus} (DS 1002); Florence, \textit{Laetentur caeli} (DS 1306); CCC 1035 \\
Inequality of punishments proportioned to the gravity of sin & \textit{De fide} & Matt 11:22--24; Lyons II (DS 858); Florence, \textit{Laetentur caeli} (DS 1306) \\
Bodily resurrection of the wicked to appear before Christ's tribunal & \textit{De fide} & Dan 12:2; John 5:28--29; \textit{Quicumque} (DS 76); Lateran IV (DS 801); \textit{Benedictus Deus} (DS 1002) \\
God predestines no one to hell; damnation results from free, unrepented mortal sin persisted in to the end & \textit{De fide} / Catholic doctrine & Ezech 33:11; 1 Tim 2:4; 2 Pet 3:9; CCC 1033, 1037, 1861; John Paul II (28 July 1999) \\
Imperfect contrition (attrition) conceived from fear of hell, excluding the will to sin with hope of pardon, is good, a gift of God, and an impulse of the Holy Ghost & \textit{De fide} & Council of Trent, Sess.\ XIV, ch.\ 4 and can.\ 5 (DS 1678, 1705); St.\ Augustine, \textit{Hom.\ on 1 John} 9.4--6; St.\ Thomas, \textit{ST} II-II, q.\ 19, aa.\ 4--7 \\
Primacy of the pain of loss (\textit{poena damni}) over the pain of sense (\textit{poena sensus}) & \textit{Sententia communis} / Ordinary Magisterium & St.\ Irenaeus, \textit{Adv.\ Haer.} V.27.2; St.\ John Chrysostom, \textit{Hom.\ on Matt.} 23.9; St.\ Gregory Nazianzen, \textit{Or.} 16.9; St.\ Augustine, \textit{Enchir.} 112; St.\ Thomas, \textit{ST} I-II, q.\ 87, a.\ 4; \textit{Roman Catechism} I.7; CCC 1035 \\
Fixed obstinacy of the damned will after death, excluding direct repentance & \textit{Sententia communis} & St.\ John of Damascus, \textit{De fide orth.} II.4; \textit{Supplement}, q.\ 98, aa.\ 1--2; St.\ Alphonsus, \textit{Prep.\ for Death} XXVI--XXVII; CCC 1033, 1861 \\
Venial sin alone does not incur eternal punishment & Catholic doctrine & St.\ Gregory the Great, \textit{Dial.} IV.39; St.\ Thomas, \textit{ST} I-II, q.\ 87, a.\ 5 \\
Physical nature of hell's fire (corporeal fire of the same species as terrestrial fire versus non-material fire known to God) & Open theological question & St.\ Augustine, \textit{City of God} XXI.10, and \textit{Supplement}, q.\ 97, aa.\ 5--6, versus St.\ John of Damascus, \textit{De fide orth.} IV.27 \\
Nature of the undying worm (corporeal versus spiritual remorse of conscience) & Open theological question & St.\ Augustine, \textit{City of God} XXI.9 (expressly leaving either open); \textit{Supplement}, q.\ 97, a.\ 2 (preferring spiritual remorse) \\
Relative proportion of the saved and the lost & Open theological question & St.\ John Chrysostom, \textit{Hom.\ on Acts} 24, and St.\ Thomas, \textit{ST} I, q.\ 23, a.\ 7, ad 3, versus Pope Benedict XVI, \textit{Spe salvi} 46; Roman Missal Secret \textit{pro vivis et defunctis} \\
Final damnation of any named individual human being & Undetermined by dogmatic definition & CCC 1861; St.\ Augustine, \textit{City of God} XXI.24; Judas noted in Matt 26:24, John 17:12, and the 1962 Holy Thursday and Good Friday Collect \\
Visions of hell reported by St.\ Teresa of Avila (\textit{Life} ch.\ 32) and at Fatima (13 July 1917) & Approved private testimony & Ecclesiastical recognition of authenticity and spiritual fruit; adds nothing to public Revelation and does not bind the assent of divine faith \\
\end{threestable}
